% Sections 6.1-6.7, from lectures/L06/appendix/a_counters_and_shift_registers.md (A.1-A.7).

\section{From register to counter}
\label{c6:sec:counter}

A register (\chapref{c3:ch}) is a bank of D flip-flops storing whatever value is fed in at the next
rising edge. A \textbf{counter} is the same idea with a twist: instead of a new value from outside,
you feed the register its \emph{own current value plus one}. At every edge it captures
\code{current_value + 1}, so the stored number increases once per clock cycle with no further
input.

\subsection*{Build it by hand first}
As with every circuit in this book: build it before you write it. A 4-bit counter is three
elements, and watching it run is the whole point:
\begin{itemize}
  \item A \textbf{4-bit register}: four D flip-flops sharing a clock, exactly the register from
    \secref{c3:sec:register}.
  \item An \textbf{adder} and a \textbf{constant \code{1}}: the register's four outputs to one of
    the adder's inputs, the constant to the other.
  \item The adder's sum wired back to the register's four \code{D} inputs. That feedback is the
    whole circuit: the register holds a number, the adder computes that number plus one, the next
    edge stores it. It is \secref{c3:sec:seq}'s feedback loop again, with an adder in the way
    instead of a plain wire.
\end{itemize}

\lecfig[0.7]{lectures/L06/appendix/images/counter_circuit.png}{A 4-bit counter: a register whose own
output, plus one, is wired back to its input.}{c6:fig:countercircuit}

Then confirm, one clock edge at a time, with a period of \code{1000 ms}:
\begin{itemize}
  \item Watch the outputs step \code{0000, 0001, 0010, 0011, ...}, one count per rising edge.
  \item \textbf{Let it run past \code{1111}.} It returns to \code{0000} by itself, and nothing in
    your circuit tells it to: there is no comparison, no clearing logic, no ``if the count is at the
    maximum'' anywhere. The adder produced \code{10000}, the fifth bit had nowhere to go, and the
    four bits that survived are \code{0000}. That is \textbf{overflow}, and in a counter it is not a
    bug to work around but the property the rest of this chapter and the whole of \chapref{c7:ch}
    rest on.
  \item \textbf{Keep an eye on \code{carry_out} on the adder while it happens.} That is the fifth
    bit, and the drawing shows exactly where it goes: high during the single edge where the count
    wraps, connected to nothing, out of the adder and no further. A wider counter would feed it into
    a fifth stage; here it is thrown away, and throwing it away \emph{is} the wraparound. So ``the
    count wraps because the bits ran out'' is not a figure of speech. You can point at the bit that
    ran out.
\end{itemize}

Keep the project. \Chapref{c7:ch} starts from exactly this circuit and adds two things to it, and
reopening it beats redrawing it.

\subsection*{The same counter in VHDL}
One signal and one process, exactly like \chapref{c3:ch}'s register, except that the right-hand side
refers to the signal itself:

\begin{vhdlcode}
signal counter: natural range 0 to 15;
...
process(clock, reset_s2_n) is
begin
    if (reset_s2_n = '0') then
        counter <= 0;
    elsif (rising_edge(clock)) then
        counter <= counter + 1;
    end if;
end process;
\end{vhdlcode}

\code{natural} is a predefined integer subtype in VHDL: a non-negative integer, unlike the
\code{std_logic} types which model bits directly. VHDL's arithmetic operators work on it the way
ordinary integer maths does, which is why \code{counter + 1} needs no extra machinery. Synthesis
turns it into a register of flip-flops all the same, exactly as a \code{std_logic_vector} would
become; the difference is purely in how you write and reason about the value.

The range constraint \code{0 to 15} is not just documentation. It tells the synthesis tool that this
counter needs four bits, so counting past \code{15} falls back to \code{0} for free, purely as a
consequence of the bits running out, with no explicit ``if counter = 15'' logic in either the
drawing or the VHDL.

The free wraparound has one condition worth saying out loud: it happens only when the upper bound is
$2^N - 1$. \code{0 to 15} wraps for free; \code{0 to 9} does not, since \code{10} is not a
power-of-two boundary, and a modulo-10 counter needs an explicit comparison. That is exactly what
Exercise~\ref{c6:ex:counter} asks you to write.

\begin{warning}
\textbf{One place where the simulator disagrees with the hardware.} The free wraparound is a
property of the \emph{synthesised} circuit, and the drawing shows that it is real. A VHDL
\emph{simulator} is stricter than a wire:
\begin{itemize}
  \item \code{natural range 0 to 15} is range-checked, so GHDL treats $15 + 1 = 16$ as an error and
    aborts instead of wrapping.
  \item The synthesised hardware has no such notion. There is no sixteenth value for the bits to
    take, so they land on \code{0000}, exactly as your four LEDs did.
\end{itemize}
The gate network and the chip both wrap; only the simulator objects, since a range-constrained
integer is a promise you made and \code{16} breaks it. The rule that follows is not ``trust the
hardware'': it is that anything you intend to verify should say what it means. Write the comparison
out, as \code{timer} and \code{serial_rx8} do, and the drawing, the simulator and the chip all
agree.
\end{warning}

This wraparound that follows from the bit width is the counter's single most important property, and
it is what \chapref{c7:ch} builds a timer out of: a counter watching for a particular count value and
raising a flag when it arrives.

% ------------------------------------------------------------------------------------------
\section{Shift registers}
\label{c6:sec:shift}

A \textbf{shift register} is, again, N D flip-flops, wired so that each flip-flop's output feeds the
\emph{next} flip-flop's input instead of its own. At every clock edge every bit moves one position
down the chain, simultaneously:

\begin{textdrawing}
serial_in -> [D Q]-->[D Q]-->[D Q]-->[D Q] -> serial_out
              FF0     FF1     FF2     FF3
             bit 0   bit 1   bit 2   bit 3
              ^clock  ^clock  ^clock  ^clock   (all share the same clock)
\end{textdrawing}

After one edge \code{FF0}'s contents are in \code{FF1}, \code{FF1}'s in \code{FF2}, and so on; a new
bit comes into \code{FF0} from \code{serial_in}, and what was in the last flip-flop falls off the end
as \code{serial_out}. Run it for N edges and N serial bits have travelled all the way through.

\textbf{One direction, used throughout this book.} A shift register can move bits either way, and
both are equally valid; it is mixing them within one design that creates confusion. So every shift
register here follows the diagram above:
\begin{itemize}
  \item the chain maps onto a \code{std_logic_vector} with \code{FF0} as \textbf{bit 0} and
    \code{FF(N-1)} as bit \code{N-1}.
  \item serial data \textbf{enters at bit 0} and \textbf{leaves from bit \code{N-1}}.
  \item every bit moves one position towards the \textbf{most significant} bit per edge.
\end{itemize}

In VHDL that direction is a single expression, identical in every example below:

\begin{vhdlcode}
shift_reg <= shift_reg(N-2 downto 0) & serial_in;
\end{vhdlcode}

\code{shift_reg(N-2 downto 0)} is everything but the top bit, moved up one position; the top bit has
nowhere to go and falls off as \code{serial_out}; \code{serial_in} fills the vacated bit 0.

When you meet a shift register outside this book: check its direction before you assume anything.
The mirror image, \code{serial_in & shift_reg(N-1 downto 1)}, is exactly as common and at a glance
looks almost the same.

Shift registers are named after how data comes in and leaves:

\begin{booktable}{@{}p{7.6em}p{6.6em}p{6.6em}L@{}}
  \thead{Form} & \thead{Data in} & \thead{Data out} & \thead{Typical use} \\
  \midrule
  SISO & one bit per clock & one bit per clock & a plain delay line \\
  \textbf{SIPO} & one bit per clock & all N bits at once & receiving a stream of serial bits, e.g.
    an incoming byte on a single wire \\
  \textbf{PISO} & all N bits at once (loaded) & one bit per clock & sending N bits out over a single
    wire \\
  PIPO & all N bits at once & all N bits at once & in practice an ordinary register
    (\secref{c3:sec:register}) \\
\end{booktable}

This chapter covers the two practically useful ones, SIPO and PISO.

% ------------------------------------------------------------------------------------------
\section{Serial in, parallel out (SIPO)}
\label{c6:sec:sipo}

A SIPO register is what you reach for when bits arrive one at a time on a single wire and you need
the accumulated value all at once. The classic is deserialising an incoming byte, which is
fundamentally what the receive path of any serial protocol does.

Which end of the byte arrives first is each protocol's own decision. SPI and CAN send the most
significant bit first by default, which is what \secref{c6:sec:shift}'s direction assumes: the first
bit to arrive ends up in the MSB once all N have been shifted in. (Bit order in SPI is
configurable on most controllers, so check the device.) A protocol sending the least significant bit
first, like UART, shifts the mirrored way, and its idiom is the mirror image too.

The idiom is one line inside a clocked process, \secref{c6:sec:shift}'s shift expression verbatim:

\begin{vhdlcode}
signal shift_reg: std_logic_vector(7 downto 0);
...
process(clock, reset_s2_n) is
begin
    if (reset_s2_n = '0') then
        shift_reg <= (others => '0');
    elsif (rising_edge(clock)) then
        shift_reg <= shift_reg(6 downto 0) & serial_in;
    end if;
end process;
parallel_out <= shift_reg;
\end{vhdlcode}

After 8 edges the received bits have entirely replaced the register's original contents, and
\code{parallel_out} exposes all 8 at once.

% ------------------------------------------------------------------------------------------
\section{Parallel in, serial out (PISO)}
\label{c6:sec:piso}

A PISO register reverses SIPO's job: load N bits at once, then send them out one at a time. It is
what you reach for when a fixed pattern has to go out over a single wire, for example to a chain of
LEDs or a display driven by an external shift register chip.

It needs one control signal more than SIPO: something separating \emph{loading} a new parallel value
from \emph{shifting} the existing one out.

\begin{vhdlcode}
process(clock, reset_s2_n) is
begin
    if (reset_s2_n = '0') then
        shift_reg <= (others => '0');
    elsif (rising_edge(clock)) then
        if (load = '1') then
            shift_reg <= parallel_in;                          -- Parallel load.
        elsif (shift = '1') then
            shift_reg <= shift_reg(N-2 downto 0) & serial_in;  -- Shift toward the MSB.
        end if;
    end if;
end process;
serial_out <= shift_reg(N-1);
\end{vhdlcode}

Bit \code{N-1} is always the next bit to leave, so a loaded value goes out most significant bit
first: load \code{"1010"} into a 4-bit register and \code{serial_out} presents \code{1}, \code{0},
\code{1}, \code{0} over the following cycles.

Note how little separates this from \secref{c6:sec:sipo}'s SIPO. \textbf{The shift expression is
identical}, same direction, same line of VHDL. Only two things differ: PISO adds a \code{load}
branch, and it taps \code{shift_reg(N-1)} as a serial output where SIPO exposes the whole vector.
``SIPO'' and ``PISO'' name how you \emph{wire} a shift register, not two kinds of hardware, which is
why \secref{c6:sec:shift} could describe the shifting once before either name came up.

Whether a design needs a genuinely new \code{serial_in} at every shift, or is content to recirculate
what falls off the end, is decided by how it is wired. That recirculation is exactly the trick the
walking LED example in \chapref{c7:ch} uses, once it has a timer clocking the shifting.

% ------------------------------------------------------------------------------------------
\section{Worked example: an 8-bit serial receiver}
\label{c6:sec:rx8}

The chapter's two halves meet in one module. A shift register can receive serial bits forever but
cannot say \emph{when} a complete byte has arrived, since it has no notion of ``how many bits so
far''. That is a counting problem, and \secref{c6:sec:counter} already solved it.
\code{serial_rx8} combines the two, plus a one-cycle \code{data_ready} pulse when the eighth bit
lands.

This is no toy: it is the front half of every serial receiver there is. The receive paths of SPI,
CAN and UART all start with a shift register and a bit counter, and differ only in what decides when
a bit is valid.

\begin{booktable}{@{}p{6.6em}p{3.4em}p{9em}L@{}}
  \thead{Port} & \thead{Dir.} & \thead{Type} & \thead{Meaning} \\
  \midrule
  \code{clock} & in & \code{std_logic} & System clock. \\
  \code{reset_s2_n} & in & \code{std_logic} & Active-low, \textbf{already synchronised} reset;
    clears the register and the bit count. Asynchronous, so it clears them with no clock edge
    involved. \\
  \code{shift_enable} & in & \code{std_logic} & High for exactly one clock cycle per incoming
    bit. \\
  \code{serial_in} & in & \code{std_logic} & The incoming bit, sampled at every enabled rising
    edge. \\
  \code{data_out} & out & \code{std_logic_vector(7 downto 0)} & The byte received so far, most
    significant bit first. \\
  \code{data_ready} & out & \code{std_logic} & One-cycle pulse when the eighth bit lands. \\
\end{booktable}

\begin{aside}
\textbf{The reset port is the synchronised one, and \code{_s2} says so.} It is called
\code{reset_s2_n} rather than \code{reset_n} because what belongs on it is the output of a
\code{reset_sync} (\secref{c4:sec:resetsync}), the module you wrote in
Exercise~\ref{c4:ex:resetsync}: assert asynchronously, release synchronously. \textbf{Instantiate
one in the design wrapping this receiver and connect its \code{reset_s2_n} here}, which is exactly
what \secref{c6:sec:shiftdemo}'s board wrapper does. Wire a push button straight into this port
instead and you get the race \chapref{c4:ch} exists to prevent: the shift register and the bit
counter are a dozen flip-flops, and an asynchronously \emph{released} reset lets them leave reset at
different edges, so the receiver can start counting a byte one cycle before it starts shifting one
in.

\code{counter}, \code{sipo8} and \code{piso8} in the exercises take \code{reset_s2_n} for the same
reason. None of them is a design you put on a board by itself: a counter, a shift register and a
receiver are parts, and a part is entitled to assume that the design around it has handled the
synchronisation. So this whole chapter takes the synchronised reset signal, and only the wrapper
modules in \secref{c6:sec:counterdemo} and \ref{c6:sec:shiftdemo}, which own the pins, ever see a raw
\code{reset_n}.

A \textbf{top level} is named the other way round: it takes \code{reset_n} and instantiates
\code{reset_sync} itself, since it is the one holding the pin. \code{blinker} and
\code{walking_led} in \chapref{c7:ch}, and both capstones in \chapref{c8:ch}, are written that way.
Get into the habit of reading \code{_s2} as a requirement on the caller, because from here on most
of what you write is a subcomponent.
\end{aside}

\vhdlfile{lectures/L06/serial_rx8/serial_rx8.vhd}

At every rising edge where \code{shift_enable = '1'}, \code{serial_in} is shifted into bit 0 with
\secref{c6:sec:sipo}'s idiom unchanged, so the first bit to arrive has travelled up to bit 7 by the
time the eighth lands, and an internal bit counter of type \code{natural range 0 to 7} increments.
At the edge where that counter already stands at \code{7}, \code{data_ready} pulses high for one
cycle and the counter clears back to \code{0}, ready for the next byte. \code{data_out} is driven by
a concurrent assignment from the shift register, so the complete byte is on it during the cycle
where \code{data_ready} is high.

Note that the counter is \emph{cleared} here rather than left to wrap: \code{0 to 7} is a
power-of-two range so it would wrap for free, but writing it out keeps the simulator in agreement
with the hardware, per \secref{c6:sec:counter}.

Two details are worth dwelling on, since both are patterns you have met more than once:
\begin{itemize}
  \item \textbf{\code{data_ready} is a pulse, not a level.} It is driven low unconditionally at the
    top of the clocked branch and raised only at the edge completing a byte: the same one-cycle
    shape as the edge detector in \secref{c3:sec:edge}. A consumer that misses it has missed the
    byte, which is exactly why it is paired with a \code{data_out} that \emph{holds} its value
    afterwards. A real peripheral would put a holding register or a small FIFO here so that a slow
    consumer cannot lose data.
  \item \textbf{\code{shift_enable} gates the bit counter, not just the shifting.} Both advance
    together or neither does, so a gap in the incoming stream pauses the receiver mid-byte instead
    of corrupting it, and it resumes on exactly the bit it was waiting for.
\end{itemize}

\paragraph{Checking it}
The example comes with a \file{serial_rx8_tb.vhd}, which shifts in two bytes with a pause between
them and checks that \code{data_ready} pulses once per byte, never early:

\begin{shell}
cd lectures/L06/serial_rx8
ghdl -a --std=93 serial_rx8.vhd serial_rx8_tb.vhd
ghdl -e --std=93 serial_rx8_tb
ghdl -r --std=93 serial_rx8_tb --assert-level=error --stop-time=10ms
\end{shell}

Clocked at \code{50 MHz} and continuously enabled, this receiver swallows a byte in 160~ns, which no
LED can show you. That is why the testbench matters here in a way it did not for a blinking LED, and
that is how most of real digital design goes. But it is no reason to skip the board, because
\code{shift_enable} is precisely the way out: see \secref{c6:sec:shiftdemo}.

% ------------------------------------------------------------------------------------------
\section{On the board: watching a counter wrap}
\label{c6:sec:counterdemo}

The counter is demonstrated on the DE0-CV during the session, built the same way as every design
since \chapref{c1:ch}: a Quartus Prime Lite project against \code{5CEBA4F23C7N}, the \file{.vhd}
file added, compile, assign every port a physical pin in the Pin Planner, then recompile and program
the board.

One thing has to change first, and it is worth seeing why. Clocked at \code{50 MHz}, a 4-bit counter
runs through all sixteen states and wraps every \code{320 ns}, about three million times a second. An
LED driven straight from it does not blink, it simply glows at half brightness. To see a wraparound
you have to watch a much slower bit: count in a wider register, say 28 bits, and assign the
\textbf{top} four bits to the LEDs.

The demonstration is therefore a wrapper around two modules rather than the exercise's entity on its
own: a \code{reset_sync} taking the raw \code{reset_n} from the pin, and the widened counter taking
its \code{reset_s2_n} from there. The counter's port is called \code{reset_s2_n} precisely so that
this connection is the only one that type-checks, and \secref{c6:sec:shiftdemo} repeats the pattern
with a third module.

\begin{itemize}
  \item \code{clock} to the 50 MHz oscillator, \code{reset_n} to a push button (into
    \code{reset_sync}, never straight into the counter), \code{count(27 downto 24)} to four LEDs.
  \item Bit 24 changes once every $2^{24}$ clock cycles, about a third of a second, so the four LEDs
    count visibly from \code{0000} up to \code{1111} and then, with nothing in the circuit doing it,
    back to \code{0000}. A full round takes about five seconds.
\end{itemize}

Dividing a clock by watching a high bit is the crudest timer imaginable, and its limitation is worth
naming now: the period is locked to a power of two, and the only way to change it is to pick a
different bit. \Chapref{c7:ch} replaces it with a counter compared against a target value you
choose, which is all a timer really is.

% ------------------------------------------------------------------------------------------
\section{On the board: shifting in a byte by hand}
\label{c6:sec:shiftdemo}

\Secref{c6:sec:rx8} said that a byte arrives too fast to follow. The solution is not to slow the
clock down, which would be the software instinct and is the wrong move in hardware: the clock is the
one thing in a synchronous design you do not fiddle with. \code{serial_rx8} already has the right
control, and it is \code{shift_enable}, which ``gates the bit counter, not just the shifting''. Hold
it low and the receiver stops dead, mid-byte, for as long as you like. Raise it for exactly one
cycle and exactly one bit goes in.

So drive \code{shift_enable} from a \textbf{button} instead of holding it high. One press, one bit.

The demonstration is a small wrapper around three modules you already have, and no new logic beyond
a single flip-flop:
\begin{itemize}
  \item \code{reset_sync} (\secref{c4:sec:resetsync}), turning the raw \code{reset_n} into
    \code{reset_s2_n} for everything else.
  \item \code{button_sync} (\secref{c4:sec:buttonsync}) with \code{COUNT = 1}, turning a press into
    the one-cycle pulse \code{shift_enable} wants. This is no optional garnish: a raw button
    connected to \code{shift_enable} would shift in a random handful of bits per press, one per
    clock cycle for as long as the contacts bounce.
  \item \code{serial_rx8} itself, unchanged, with its \code{reset_s2_n} taken from
    \code{reset_sync} above rather than from the button. That is the whole reason the port is called
    \code{reset_s2_n}: the wrapper owns the raw \code{reset_n}, and the receiver only ever gets to
    see the synchronised one.
  \item A flip-flop latching \code{data_ready}, set when the pulse arrives and cleared by reset.
    Without it there is nothing to see: \code{data_ready} is high for 20~ns, and no LED can show you
    that.
\end{itemize}

On the board: \code{clock} to the 50 MHz oscillator, \code{reset_n} to a push button,
\code{serial_in} to a switch (that is the bit you are about to send), the shift button to a second
push button, \code{data_out(7 downto 0)} to eight LEDs, and the latched \code{data_ready} to a
ninth.

Now set the switch, press the button, and watch one bit come in at the rightmost LED and the whole
pattern step left. Send \code{10110010} one bit at a time and the byte builds up across the LEDs in
front of you, most significant bit first, exactly as \secref{c6:sec:sipo}'s idiom says it must. At
the eighth press the ninth LED lights and the bit counter starts over.

Three things this makes visible that the testbench can only assert:
\begin{itemize}
  \item \textbf{The bit order.} The first bit you send ends up in bit 7. Reading that off a waveform
    is a chore; watching it travel across a row of LEDs is not.
  \item \textbf{What \code{shift_enable} is for.} Let go of the button for a minute mid-byte and
    nothing moves, nothing is lost, and the eighth bit still completes the byte. That is the
    ``pauses mid-byte instead of corrupting it'' property, demonstrated rather than claimed.
  \item \textbf{Pulse versus level, once more.} The reason a flip-flop has to be added for
    \code{data_ready} is that whole distinction, and it is the same reason the edge detector in
    \secref{c3:sec:edge} existed.
\end{itemize}
